\documentclass[11pt]{article}
\usepackage[margin=1in]{geometry}
\usepackage{amsmath,amssymb,amsthm}
\usepackage{array}
\usepackage{parskip}
\usepackage{booktabs}
\setlength{\emergencystretch}{3em}

\newtheorem{theorem}{Theorem}
\newtheorem{lemma}[theorem]{Lemma}
\newtheorem{corollary}[theorem]{Corollary}
\newtheorem{proposition}[theorem]{Proposition}
\theoremstyle{definition}
\newtheorem{definition}[theorem]{Definition}
\theoremstyle{remark}
\newtheorem{remark}[theorem]{Remark}

\newcommand{\B}{\mathcal{B}}
\newcommand{\FB}{\mathrm{FB}}

\title{Disproof of conjecture \texttt{00000002051}}
\author{earthking11}
\date{2026-09-13}

\begin{document}
\maketitle

\section*{Verdict: FALSE}

We disprove conjecture \texttt{00000002051} as stated. The conjecture claims that
the fine spectrum of the semigroup variety $\langle x^2 = x\rangle$ is the closed
form $f(n) = 2^n C_n$. The variety $\langle x^2 = x\rangle$ is the variety of
\emph{bands} (idempotent semigroups), and the fine spectrum is by definition the
cardinality of the $n$-generated free algebra. The formula fails already at
$n = 1$: the free band on one generator is a singleton, so $f(1) = 1$, whereas
the formula gives $2^1 C_1 = 2$. An independent second failure at $n = 2$
($f(2) = 6 \ne 8$) shows the claim cannot be repaired by restricting or shifting
the range of $n$.

\section{The conjecture and its terms}

Quoted verbatim from \texttt{conjectures/00000002051.md}:

\begin{quote}
\textbf{English.} Definition: The fine spectrum $f_V(n)$ of a variety $V$ is the
cardinality of its $n$-generated free algebra. Conjecture: The fine spectrum of
the semigroup variety $\langle x^2 = x\rangle$ is the explicit closed form
$f(n) = 2^{n}\cdot C_n$ ($C_n$ the Catalan numbers); the spectrum is controlled
by the connected-component count of band semigroups.
\end{quote}

The file fixes the meaning of its two key terms itself:

\begin{definition}[fine spectrum]
$f_V(n) = |F_V(n)|$, the number of elements of the free algebra of the variety
$V$ on $n$ generators.
\end{definition}

\begin{definition}[the variety $\langle x^2 = x\rangle$]
A semigroup in which $x^2 = x$ for every element $x$ is an \emph{idempotent
semigroup}, i.e.\ a \emph{band}. Thus $F_V(n)$ is the free band on $n$
generators, and $f(n)$ is its cardinality. The file's own concluding clause
(``band semigroups'') confirms this reading.
\end{definition}

The natural range is $n \ge 1$. The claim is a single closed form asserted for
this range; a single value of $n$ at which the two sides differ refutes it.

\section{Failure at $n = 1$}

\begin{theorem}\label{thm:one}
The free band on one generator has exactly one element; that is, $f(1) = 1$.
The conjecture's formula gives $2^1 C_1 = 2 \ne 1$.
\end{theorem}

\begin{proof}
Let $B$ be the free band on the single generator $x$, and let $S \subseteq B$ be
the subsemigroup generated by $x$. Every element of $S$ is a finite product
$x x \cdots x$. Since $B$ is a band, $x \cdot x = x$; multiplying a product of
$k$ copies of $x$ by one more copy gives $x^k \cdot x = x^{k+1}$, and induction
on $k$ using $x \cdot x = x$ shows $x^k = x$ for every $k \ge 1$. Hence every
element of $S$ is $x$, so $S = \{x\}$ is a singleton and $B = S$. Therefore
$|B| = f(1) = 1$.

More generally, in \emph{any} band $B$ and for any element $x \in B$, the
subsemigroup generated by $x$ is the singleton $\{x\}$: the set $\{x\}$ is
closed under multiplication precisely because $x \cdot x = x$. This universal
statement is what makes the refutation independent of any structural theory of
bands. The free band on one generator is unique up to isomorphism and has one
element, so
\[
f(1) = 1 \;\ne\; 2 = 2^1 C_1 .
\]
\end{proof}

\begin{corollary}
Conjecture \texttt{00000002051} is false. The closed form $f(n) = 2^n C_n$
already disagrees with the fine spectrum at the smallest admissible value
$n = 1$.
\end{corollary}

\begin{remark}
The argument is a single line: $2^1 C_1 = 2$ but the free band on one generator
is $\{x\}$. No generating set of size $1$ can produce two distinct elements,
because the only word one can form is $x$ itself.
\end{remark}

\section{Failure at $n = 2$}

The second failure is independent of Theorem~\ref{thm:one}.

\begin{theorem}\label{thm:two}
The free band on two generators has exactly six elements
\[
\{\,a,\; b,\; ab,\; ba,\; aba,\; bab\,\}.
\]
Hence $f(2) = 6$, whereas the conjecture's formula gives $2^2 C_2 = 4 \cdot 2 =
8$.
\end{theorem}

\begin{proof}
Let $a, b$ be the free generators. The free band is the quotient of the free
semigroup $\{a,b\}^+$ by the congruence generated by $u \cdot u = u$ for all
words $u$ (idempotency of every element). We first observe that \emph{every
word of length $\ge 4$ over a two-letter alphabet contains a square factor}, a
factor $ss$ consisting of two equal adjacent blocks. Indeed it suffices to check
the $16$ words of length $4$ over $\{a,b\}$; each contains a square factor
(among them $aa$, $bb$, $abab$, $abba$, $baab$, $baba$), and every word of
length $\ge 4$ has its first four letters as a factor. Consequently any word
with a factor $ss$ reduces, by $s \cdot s = s$, to the strictly shorter word
with one copy of $s$ removed; iterating, every word reduces to a
\emph{square-free} word
of length $\le 3$. The square-free words of length $\le 3$ over $\{a,b\}$ are
exactly the six displayed words (length $1$: $a, b$; length $2$: $ab, ba$;
length $3$: $aba, bab$; the words $aa$, $bb$, $aab$, $abb$, $aaa$, $bbb$, \dots{}
all contain a square). Therefore $f(2) \le 6$.

For the reverse inequality, the induced multiplication on the six reduced words
(taking reduced concatenation) is the table below; the entry in row $u$ and
column $v$ is the reduced form of the concatenation $uv$.
\end{proof}

\begin{center}
\begin{tabular}{c|cccccc}
\toprule
$\cdot$ & $a$ & $b$ & $ab$ & $ba$ & $aba$ & $bab$ \\
\midrule
$a$   & $a$   & $ab$  & $ab$  & $aba$ & $aba$ & $ab$  \\
$b$   & $ba$  & $b$   & $bab$ & $ba$  & $ba$  & $bab$ \\
$ab$  & $aba$ & $ab$  & $ab$  & $aba$ & $aba$ & $ab$  \\
$ba$  & $ba$  & $bab$ & $bab$ & $ba$  & $ba$  & $bab$ \\
$aba$ & $aba$ & $ab$  & $ab$  & $aba$ & $aba$ & $ab$  \\
$bab$ & $ba$  & $bab$ & $bab$ & $ba$  & $ba$  & $bab$ \\
\bottomrule
\end{tabular}
\end{center}

\begin{proposition}\label{prop:checks}
The table is closed, associative, and idempotent.
\end{proposition}

\begin{proof}
Closure is immediate from the table (all entries lie in the six-word set).
Associativity is the assertion $u(vw) = (uv)w$ for all $6^3 = 216$ triples;
idempotence is $u^2 = u$ for the six diagonal entries. Both were verified by
exhaustive computation in \texttt{reproduce.py}, which reduces all words over
$\{a,b\}$ up to length $8$ and checks the table entry by entry. The six
diagonal entries are $a, b, ab, ba, aba, bab$, so the checks pass.
\end{proof}

\begin{corollary}
The six-element table is a band. It is generated by the two elements $a$ and
$b$, so by the universal property of the free band it is a quotient of the free
band on two generators, giving $f(2) \ge 6$. Combined with $f(2) \le 6$ from
the reduction argument, $f(2) = 6$. The conjecture's value is $8$, so it fails
at $n = 2$ as well.
\end{corollary}

\begin{remark}
Six is a classical value: the free band on two generators has six elements
(Howie, \emph{Fundamentals of Semigroup Theory}, Oxford University Press, 1995,
p.~123). It is the second term of OEIS A030449, whose terms are
$1, 6, 159, 332380, \dots$.
\end{remark}

\section{True free-band values versus the conjectured formula}

The number of elements of the free band on $n$ generators is OEIS
\textbf{A030449} (Howie 1995, p.~123). The values differ from $2^n C_n$ for
every $n \ge 1$, and the discrepancy grows:

\begin{center}
\begin{tabular}{r r r l}
\toprule
$n$ & free band $f(n)$ (A030449) & $2^n C_n$ & equal? \\
\midrule
$1$ & $1$      & $2$   & no \\
$2$ & $6$      & $8$   & no \\
$3$ & $159$    & $40$  & no \\
$4$ & $332380$ & $224$ & no \\
\bottomrule
\end{tabular}
\end{center}

\noindent
Here $C_1, C_2, C_3, C_4 = 1, 2, 5, 14$. The formula overcounts at $n = 1, 2$
and drastically undercounts from $n = 3$ on: $2^n C_n$ is a genuine sequence
$2, 8, 40, 224, \dots$, but it is not the band fine spectrum.

\section{Alternative readings}

The refutation is robust under the readings a defender of the conjecture might
propose.

\begin{enumerate}
\item \emph{``$\langle x^2 = x\rangle$ means a different variety or a
free object $\langle x \mid x^2 = x\rangle$.''} A semigroup with $x^2 = x$ is an
idempotent semigroup, i.e.\ a band, and the file says ``band semigroups''. The
presentation $\langle x \mid x^2 = x\rangle$ has a single generator $x$; its
free object on one generator is still $\{x\}$. Both readings agree and both give
$f(1) = 1 \ne 2$.

\item \emph{``$f(n)$ counts irreducible or maximal words rather than algebra
elements.''} The definition says ``cardinality of its $n$-generated free
algebra''. Even under this alternative reading the counts for $n = 1, 2$ are the
irreducible-word counts $1$ and $6$, hence still differ from $2$ and $8$.

\item \emph{``$n$ starts at $2$.''} Then $f(2) = 6 \ne 8 = 2^2 C_2$; the claim
is still false. Shifting further does not help: $f(3) = 159 \ne 40$ and
$f(4) = 332380 \ne 224$, and the two sequences never agree for $n \ge 1$.

\item \emph{``Adjoin an identity to the free band.''} Under this non-standard
reading $f(1) = 2 = 2^1 C_1$ happens to match, but $f(2) = 6 + 1 = 7 \ne 8 =
2^2 C_2$, so the claim still fails. (The conjecture speaks of a \emph{semigroup}
variety, where no identity is adjoined, so this is not the intended reading in
any case.)

\item \emph{$n = 0$.} The free semigroup, hence free band, on the empty set is
empty, so $f(0) = 0$, whereas the formula gives $2^0 C_0 = 1$. If one instead
declares the $0$-generated free band to be trivial, then $n = 0$ happens to
match, but the failures at $n = 1$ and $n = 2$ are unaffected.
\end{enumerate}

No reading we could construct saves the claim. The refutation is unconditional,
and the cheapest formalisable witness is $n = 1$: since the formula is asserted
as a closed form for all $n \ge 1$, the single value $f(1) = 1 \ne 2 = 2^1 C_1$
already suffices.

\section{Reproducibility and formalisation}

The computations of Sections~4 and~5 are reproduced by
\texttt{reproduce.py}, which uses the Python~3 standard library only. It
enumerates all words over $\{x\}$ and over $\{a,b\}$ up to length $8$, reduces
them by $u \cdot u \to u$ to a fixed point (collapse the longest square factor
first), and asserts

\begin{itemize}
\item for $n = 1$: exactly one reduced class, $\{x\}$, so $f(1) = 1 \ne 2 =
2^1 C_1$;
\item for $n = 2$: exactly six reduced classes $\{a,b,ab,ba,aba,bab\}$, with a
multiplication table that is closed, associative (all $216$ triples) and
idempotent, and generated by $\{a,b\}$, so $f(2) = 6 \ne 8 = 2^2 C_2$;
\item every word of length $4$ through $12$ over two letters has a square
factor (the mechanism that makes the six words complete);
\item the literature values $1, 6, 159, 332380$ (A030449) differ from
$2^n C_n = 2, 8, 40, 224$ for $n = 1, 2, 3, 4$;
\item the alternative readings of Section~6 also fail.
\end{itemize}

\noindent
It prints \texttt{PASS} and exits $0$ exactly when every check holds, and prints
\texttt{FAIL} with a non-zero exit status otherwise.

The accompanying Lean~4 project under \texttt{lean4/} formalises the core of the
disproof in core Lean only (\texttt{import Std}, no Mathlib, no \texttt{sorry}).
It defines the structure of a \texttt{Band} (an associative, idempotent binary
operation) and proves that in \emph{any} band the subsemigroup generated by one
element is the singleton $\{x\}$; that the one-element band has the universal
property of the free band on one generator, so $f(1) = 1$; that the explicit
six-element multiplication table on \texttt{Fin 6} is associative and idempotent
by \texttt{decide} and is generated by two elements; that the Catalan numbers
satisfy $C_1 = 1$, $C_2 = 2$, $C_3 = 5$; and the arithmetic failures
$1 \ne 2^1 C_1$, $6 \ne 2^2 C_2$, $159 \ne 2^3 C_3$. The axiom audit
(\texttt{lake env lean Check.lean}) reports no \texttt{sorryAx} and no Mathlib
dependency. See \texttt{lean4/README.md} for the full statement table and the
scope note.

\section*{Files and rule-3 status}

Rule 3 requires the LaTeX source, a PDF, and a Lean~4 project.

\begin{itemize}
\item \textbf{LaTeX source} --- \texttt{main.tex}, a standalone \texttt{article}
compiling with \texttt{tectonic main.tex}.
\item \textbf{PDF} --- \texttt{build/main.pdf}, produced by \texttt{tectonic}
and placed in \texttt{build/}.
\item \textbf{Lean~4 project} --- \texttt{lean4/} with \texttt{lean-toolchain}
(\texttt{leanprover/lean4:v4.33.1}), \texttt{lakefile.toml} (library \texttt{Main},
name \texttt{tlmc2051}), \texttt{Main.lean}, \texttt{Check.lean}, and
\texttt{README.md}.
\end{itemize}

\noindent
\texttt{reproduce.py} accompanies the submission and certifies the finite
computations.

\end{document}
